\begin{tcolorbox}[colback=red!3!white,colframe=red!50!black,boxrule=0.8pt,title={Correction notes},fonttitle=\bfseries]
\textbf{Corrections from the calculation review (30~Sept.~2026).} The following places are corrected or annotated in the text; each carries a dated note there with the previous value:
\begin{itemize}
\item CMB temperature: $\frac{16}{9}\xi^2\cdot7500$: $2.35\times10^{-4} \to 2.37\times10^{-4}$; in SI $2.725$~K $\to 2.751$~K (see R70).
\item $\rho_{\text{CMB}} = 4.87\times10^{41}$: cannot be reproduced, annotated.
\item Casimir: energy density $\pi^2/(240\,d^4) \to \pi^2/(720\,d^4)$ ($240$ is the pressure); ratio to $\rho_{\text{CMB}}$ $308 \to 102.8$, direct calculation $312 \to 104$ (calculated, not a measured value), also in the equation and table further below.
\item Resonance frequency: $\nu_\xi = 1/L_\xi = 10^4$~Hz $\to c/L_\xi \approx 3\times10^{12}$~Hz.
\item $G_{\text{T0}}$: already corrected in the text on 29~Sept.~2026.
\end{itemize}
\end{tcolorbox}

\section*{Abstract}
		The T0-theory examines how far cosmic phenomena can be traced back to a single universal constant $\xi = \frac{4}{3} \times 10^{-4}$. This document presents the fundamental relationships between the gravitational constant, cosmic microwave background radiation (CMB), Casimir effect and cosmic structures within the framework of a static, eternally existing universe. All derivations are performed in natural units ($\hbar = c = k_B = 1$) and respect the time-energy duality as a fundamental principle of quantum mechanics.

	\medskip\noindent\textbf{Note on versions.} The German version of this document is more extensive than this English one. Earlier revisions shortened or changed parts of the English version. Where the two differ, the more extensive German version applies.


	\medskip\noindent\textbf{Status markers.} Statements marked \textbf{[S]} are \emph{speculative}: a hint or conjecture that has not yet been derived or numerically checked in the corpus. This document deals almost exclusively with the cosmic sector (static universe, CMB, redshift, dark energy). The current corpus deliberately leaves this sector aside, because the Standard Model does not derive it either (P39); the cosmological conclusions of this document are therefore to be read as \textbf{[S]}.

	

	\section{Introduction: The Universal $\xi$-Constant}
	
\subsection{Foundations of T0 Theory}

\begin{important}
	T0 theory is based on the universal dimensionless constant $\xi = \frac{4}{3} \times 10^{-4}$, to which it traces phenomena from the subatomic to the cosmic scale.
\end{important}

T0 theory is built on a single fundamental constant. This constant forms the basis for all physical calculations and predictions of the theory:

\begin{equation}
	\xi = \frac{4}{3} \times 10^{-4} = 1.333333... \times 10^{-4}
\end{equation}

This dimensionless constant connects quantum and gravitational phenomena and is intended as the basis of a unified description of the fundamental interactions.

\begin{tcolorbox}[colback=yellow!10!white,colframe=yellow!50!black,title=Note on Derivation]
	For the detailed derivation and physical justification of this fundamental constant, see the document ``Parameter Derivation'' (available at: ).
\end{tcolorbox}
	
	\subsection{Time-Energy Duality as Foundation}
	
	\begin{revolutionary}[Thesis]
		\textbf{[S]} From Heisenberg's uncertainty relation $\Delta E \times \Delta t \geq \hbar/2 = 1/2$ (natural units) this document concludes that a classical Big Bang singularity with infinite density is not physically realized and is replaced by a tiny but finite core with minimal length scale $L_0$ (from $\xi$).
	\end{revolutionary}
	
	Heisenberg's uncertainty relation between energy and time represents the fundamental principle of T0-theory:
	
	\begin{equation}
		\Delta E \times \Delta t \geq \frac{1}{2} \quad \text{(natural units)}
	\end{equation}
	
	In the reading of this document, this relation has far-reaching cosmological consequences \textbf{[S]}:
	\begin{itemize}
		\item A temporal beginning (Big Bang) would mean $\Delta t$ = finite
		\item This leads to $\Delta E \to \infty$ - physically inconsistent
		\item Therefore the universe must have existed eternally: $\Delta t = \infty$
		\item The universe is static, without expanding space
	\end{itemize}

	\section{Cosmic Microwave Background (CMB)}
	
	\subsection{CMB without Big Bang: $\xi$-Field Mechanisms}
	
	\begin{revolutionary}[Thesis]
		\textbf{[S]} If time-energy duality excludes a Big Bang, as concluded above, the CMB must have a different origin than the z=1100 decoupling of standard cosmology.
	\end{revolutionary}
	
	T0-theory interprets the CMB \textbf{[S]} as radiation from $\xi$-field quantum fluctuations:
	
	\begin{equation}
		\frac{T_{\text{CMB}}}{E_\xi} = \frac{16}{9} \xi^2
	\end{equation}
	
	With $E_\xi = \frac{1}{\xi} = \frac{3}{4} \times 10^4$ (natural units) and $\xi = \frac{4}{3} \times 10^{-4}$ this yields:
	
	\begin{equation}
		T_{\text{CMB}} = \frac{16}{9} \xi^2 \times E_\xi = \frac{16}{9} \times 1.78 \times 10^{-8} \times 7500 = 2.37 \times 10^{-4}
	\end{equation}

\textit{(Corrected on 30~Sept.~2026: previously $2.35 \times 10^{-4}$ -- with $\xi = \frac{4}{3}\times10^{-4}$, $\frac{16}{9}\xi^2\cdot 7500$ gives $2.370\times10^{-4}$.)}
	
	\textbf{Conversion to SI units:}
	\begin{equation}
		T_{\text{CMB}} = 2.751 \text{ K}
	\end{equation}

\textit{(Corrected on 30~Sept.~2026: previously $2.725$~K -- that is the measured value; $2.370\times10^{-4}$ (with $E_\xi$ in eV) corresponds to $2.751$~K.)}
	
	The expression thus lies about 1\,\% from the observed value (cf.\ the table of dimensionless ratios: $3.13$ versus $3.16 \times 10^{-8}$).
	
	\subsection{CMB Energy Density and $\xi$-Length Scale}
	
	The CMB energy density in natural units is:
	\begin{equation}
		\rho_{\text{CMB}} = 4.87 \times 10^{41} \quad \text{(natural units, dimension } [E^4] \text{)}
	\end{equation}

\textit{(Note of 30~Sept.~2026: the value $4.87\times10^{41}$ cannot be reproduced; in natural units with energy in eV, $\rho_{\text{CMB}} = \frac{\pi^2}{15}T_{\text{CMB}}^4$ gives about $2.0\times10^{-15}$~eV$^4$.)}
	
	This energy density defines a characteristic $\xi$-length scale:
	\begin{equation}
		L_\xi = \left(\frac{\xi}{\rho_{\text{CMB}}}\right)^{1/4}
	\end{equation}
	
	\begin{formula}
		Fundamental relation of CMB energy density:
		\begin{equation}
			\rho_{\text{CMB}} = \frac{\xi}{L_\xi^4} = \frac{\frac{4}{3} \times 10^{-4}}{(L_\xi)^4}
		\end{equation}
	\end{formula}
	
	\section{Casimir Effect and $\xi$-Field Connection}
	
	\subsection{Casimir-CMB Ratio as Consistency Check}
	
	\begin{experiment}
		The ratio between Casimir energy density and CMB energy density is consistent with the characteristic $\xi$-length scale of $L_\xi = 10^{-4}$ m.
	\end{experiment}
	
	The Casimir energy density at plate separation $d = L_\xi$ is:
	\begin{equation}
		|\rho_{\text{Casimir}}| = \frac{\pi^2}{720 \times L_\xi^4} \quad \text{(natural units)}
	\end{equation}

\textit{(Corrected on 30~Sept.~2026: previously denominator $240$ -- according to the literature, $\pi^2\hbar c/(240\,d^4)$ is the Casimir pressure between ideal plates; the energy density is $\pi^2\hbar c/(720\,d^4)$.)}
	
	The ratio yields:
	\begin{equation}
		\frac{|\rho_{\text{Casimir}}|}{\rho_{\text{CMB}}} = \frac{\pi^2}{720 \xi} = \frac{\pi^2 \times 10^4}{960} \approx 102.8
	\end{equation}

\textit{(Corrected on 30~Sept.~2026: previously $\pi^2/(240\xi) = \pi^2\times10^4/320 \approx 308$, labelled as the \enquote{experimental} ratio -- with the energy density (denominator $720$) the ratio is $102.8$; it is calculated, not measured.)}
	
	\textbf{Direct calculation with SI values:}
	With $L_\xi = 10^{-4}$ m, direct calculation gives:
	\begin{align}
		|\rho_{\text{Casimir}}| &= \frac{\hbar c \pi^2}{720 \times (10^{-4})^4} = 4.33 \times 10^{-12} \text{ J/m}^3 \\
		\rho_{\text{CMB}} &= 4.17 \times 10^{-14} \text{ J/m}^3 \\
		\text{Ratio} &= \frac{4.33 \times 10^{-12}}{4.17 \times 10^{-14}} = 104
	\end{align}

\textit{(Corrected on 30~Sept.~2026: previously denominator $240$, $1.3\times10^{-11}$~J/m$^3$ and ratio $312$ -- $1.3\times10^{-11}$ is the Casimir pressure in Pa, the energy density is one third of it; the ratio is calculated, not a measured value, only $\rho_{\text{CMB}}$ is measured.)}
	
	The theoretical expression (102.8) and the direct calculation with the measured $\rho_{\text{CMB}}$ (104) differ by about 1\%. \textit{(Corrected on 30~Sept.~2026: previously $308$, $312$ and $1.3\%$ -- follow-up values of the correction to the Casimir energy density.)}
	
	\subsection{$\xi$-Field as Universal Vacuum}
	
	\begin{important}
		The $\xi$-field manifests both in free CMB radiation and in geometrically constrained Casimir vacuum. In the interpretation of this document this is a hint at the physical reality of the $\xi$-field \textbf{[S]}; it is not a proof.
	\end{important}
	
	The characteristic $\xi$-length scale $L_\xi$ is the point where CMB vacuum energy density and Casimir energy density reach comparable magnitudes:
	
	\begin{align}
		\text{Free vacuum:} \quad &\rho_{\text{CMB}} = +4.87 \times 10^{41} \\
		\text{Constrained vacuum:} \quad &|\rho_{\text{Casimir}}| = \frac{\pi^2}{720 d^4}
	\end{align}

\textit{(Corrected on 30~Sept.~2026: previously denominator $240$ -- energy density instead of Casimir pressure, as above.)}
	
	\section{Cosmic Redshift without Expansion}
	
	\subsection{$\xi$-Field Energy Loss Mechanism}
	
	\begin{revolutionary}[Thesis]
		\textbf{[S]} According to this thesis, the observed cosmic redshift arises not from spatial expansion but from energy loss of photons in the omnipresent $\xi$-field. The wavelength dependence assumed here has been replaced in the current corpus by an achromatic redshift (correction K6 in Doc.~190).
	\end{revolutionary}
	
	Photons lose energy through interaction with the $\xi$-field:
	\begin{equation}
		\frac{dE}{dx} = -\xi \cdot f\left(\frac{E}{E_\xi}\right) \cdot E
	\end{equation}
	
	For the linear case $f\left(\frac{E}{E_\xi}\right) = \frac{E}{E_\xi}$ this yields:
	\begin{equation}
		\frac{dE}{dx} = -\frac{\xi E^2}{E_\xi}
	\end{equation}
	
	\subsection{Wavelength-Dependent Redshift}
	
	Integration of the energy loss equation leads to wavelength-dependent redshift:
	
	\begin{formula}
		Wavelength-dependent redshift:
		\begin{equation}
			z(\lambda_0) = \frac{\xi x}{E_\xi} \cdot \lambda_0
		\end{equation}
		where $\lambda_0$ is the emitted wavelength and $x$ is the distance traveled.
	\end{formula}
	
	This formula predicts:
	\begin{itemize}
		\item Shorter wavelength light (UV) shows greater redshift
		\item Longer wavelength light (radio) shows smaller redshift
		\item The ratio is $z_1/z_2 = \lambda_1/\lambda_2$
	\end{itemize}
	
	\begin{experiment}
		Experimental test: Comparison of radio and optical redshifts
		\begin{itemize}
			\item 21cm hydrogen line: $\nu = 1420$ MHz
			\item Optical H$\alpha$ line: $\nu = 457$ THz
			\item Predicted ratio: $z_{21\text{cm}}/z_{\text{H}\alpha} = 3.1 \times 10^{-6}$
		\end{itemize}
	\end{experiment}
	
	\section{Structure Formation in the Static $\xi$-Universe}
	
	\subsection{Continuous Structure Development}
	
	In the static T0 universe \textbf{[S]}, structure formation occurs continuously without Big Bang constraints:
	
	\begin{equation}
		\frac{d\rho}{dt} = -\nabla \cdot (\rho \mathbf{v}) + S_\xi(\rho, T, \xi)
	\end{equation}
	
	where $S_\xi$ is the $\xi$-field source term for continuous matter/energy transformation.
	
	\subsection{$\xi$-Supported Continuous Creation}
	
	The $\xi$-field enables continuous matter/energy transformation:
	
	\begin{align}
		\text{Quantum vacuum} &\xrightarrow{\xi} \text{Virtual particles} \\
		\text{Virtual particles} &\xrightarrow{\xi^2} \text{Real particles} \\
		\text{Real particles} &\xrightarrow{\xi^3} \text{Atomic nuclei} \\
		\text{Atomic nuclei} &\xrightarrow{\text{Time}} \text{Stars, galaxies}
	\end{align}
	
	Energy balance is maintained by:
	\begin{equation}
		\rho_{\text{total}} = \rho_{\text{matter}} + \rho_{\xi\text{-field}} = \text{constant}
	\end{equation}
	
	\section{Dimensionless $\xi$-Hierarchy}
	
	\subsection{Energy Scale Ratios}
	
	The $\xi$-relations used here reduce to exact mathematical ratios:
	
	\begin{longtable}{lcc}
		\caption{Dimensionless $\xi$-ratios} \\
		\toprule
		\textbf{Ratio} & \textbf{Expression} & \textbf{Value} \\
		\midrule
		\endfirsthead
		\multicolumn{3}{c}{\tablename\ \thetable{} -- Continued} \\
		\toprule
		\textbf{Ratio} & \textbf{Expression} & \textbf{Value} \\
		\midrule
		\endhead
		Temperature & $\frac{T_{\text{CMB}}}{E_\xi}$ & $3.13 \times 10^{-8}$ \\
		Theory & $\frac{16}{9}\xi^2$ & $3.16 \times 10^{-8}$ \\
		Length & $\frac{\ell_{\xi}}{L_\xi}$ & $\xi^{-1/4}$ \\
		Casimir-CMB & $\frac{|\rho_{\text{Casimir}}|}{\rho_{\text{CMB}}}$ & $\frac{\pi^2 \times 10^4}{960}$ \\
		\bottomrule
	\end{longtable}

\textit{(Corrected on 30~Sept.~2026: in the Casimir-CMB row previously $\pi^2\times10^4/320$ -- with the energy density $\pi^2\times10^4/960 \approx 102.8$.)}
\normalsize
	
	\begin{important}
		The $\xi$-relations used here consist of exact mathematical ratios:
		\begin{itemize}
			\item Fractions: $\frac{4}{3}$, $\frac{3}{4}$, $\frac{16}{9}$
			\item Powers of ten: $10^{-4}$, $10^3$, $10^4$
			\item Mathematical constants: $\pi^2$
		\end{itemize}
		No arbitrary decimal numbers occur in these expressions; they are built from $\xi$ and integer structural choices (motivated settings).
	\end{important}
	
	\section{Experimental Predictions and Tests}
	
	\subsection{Precision Measurements of Gravitational Constant}
	
	T0-theory predicts:
	\begin{equation}
		G_{\text{T0}} = 6.6745 \times 10^{-11} \text{ m}^3/(\text{kg} \cdot \text{s}^2) \quad (+0.003\,\%\ \text{relative to CODATA})
	\end{equation}
	
	\textit{(Corrected on 29~Sept.~2026: previously $6.67430000\ldots\times10^{-11}$ -- these digits are not covered by the calculation; the corpus formula $G = \xi^2/(4m_e)\cdot C_{\text{conv}}\cdot K_{\text{frak}}$ yields $6.6745\times10^{-11}$.)}
	
	This prediction can be tested by future precision measurements.
	
	\subsection{Casimir Force Anomalies}
	
	\begin{experiment}
		Prediction: Casimir force anomalies at characteristic $\xi$-length scale
		\begin{itemize}
			\item Standard Casimir law: $F \propto d^{-4}$
			\item $\xi$-field modifications at $d = L_\xi = 10^{-4}$ m
			\item Measurable deviations through $\xi$-vacuum coupling
		\end{itemize}
	\end{experiment}
	
	\subsection{Electromagnetic Resonance}
	
	Maximum $\xi$-field-photon coupling at characteristic frequency:
	\begin{equation}
		\nu_\xi = \frac{c}{L_\xi} \approx 3 \times 10^{12} \text{ Hz} = 3 \text{ THz}
	\end{equation}

\textit{(Corrected on 30~Sept.~2026: previously $\nu_\xi = 1/L_\xi = 10^{4}$~Hz $= 10$~kHz -- the conversion to SI lacked $c$; $1/L_\xi = 10^4$~m$^{-1}$ is a wavenumber, the frequency is $c/L_\xi \approx 3\times10^{12}$~Hz.)}
	
	Electromagnetic anomalies should occur at this frequency.
	
	\section{Cosmological Consequences}
	
	\subsection{Solution to Cosmological Problems}
	
	The following table shows how known problems of standard cosmology appear in the static T0 model \textbf{[S]}. The entries are interpretations in the cosmic sector, which the current corpus leaves aside (P39), not derived solutions:
	
	\begin{longtable}{p{3.5cm}p{4.5cm}p{6cm}}
		\caption{Cosmological problems: Standard vs. T0} \\
		\toprule
		\textbf{Problem} & \textbf{$\Lambda$CDM} & \textbf{T0 Interpretation} \\
		\midrule
		\endfirsthead
		\multicolumn{3}{c}{\tablename\ \thetable{} -- Continued} \\
		\toprule
		\textbf{Problem} & \textbf{$\Lambda$CDM} & \textbf{T0 Interpretation} \\
		\midrule
		\endhead
		Horizon problem & Inflation required & Infinite causal connectivity \\
		Flatness problem & Fine-tuning & Geometry stabilizes over infinite time \\
		Monopole problem & Topological defects & Defects dissipate over infinite time \\
		Lithium problem & Nucleosynthesis discrepancy & Nucleosynthesis over unlimited time \\
		Age problem & Objects older than universe & Objects can be arbitrarily old \\
		$H_0$ tension & 9\% discrepancy & No $H_0$ in static universe \\
		Dark energy & 69\% of energy density & Not required \\
		\bottomrule
	\end{longtable}
\normalsize
	
	\subsection{Parameter Reduction}
	
	\begin{revolutionary}[Parameter count]
		Parameter count: the Standard Model and $\Lambda$CDM together use 25+ parameters, T0-theory the single parameter $\xi$ (plus integer structural choices as motivated settings; SI anchors serve only for conversion).
		\begin{itemize}
			\item Standard model of particle physics: 19+ parameters
			\item $\Lambda$CDM cosmology: 6 parameters
			\item T0-theory: 1 parameter ($\xi$)
		\end{itemize}
		Arithmetically this corresponds to a 96\% reduction. The $\Lambda$CDM share concerns the cosmic sector, which the current corpus leaves aside (P39).
	\end{revolutionary}
	
	
	\begin{thebibliography}{20}
		
		\bibitem{t0_lagrangian_de}
		Pascher, Johann (2025). 
		\textit{Vereinfachte Lagrange-Dichte und Zeit-Massen-Dualität in der T0-Theorie}. 
		T0-Theory Project. 
		\url{https://jpascher.github.io/T0-Time-Mass-Duality/2/pdf/129_lagrandian-einfach_De.pdf}
		
		\bibitem{t0_lagrangian_en}
		Pascher, Johann (2025). 
		\textit{Simplified Lagrangian Density and Time-Mass Duality in T0-Theory}. 
		T0-Theory Project. 
		\url{https://jpascher.github.io/T0-Time-Mass-Duality/2/pdf/129_lagrandian-einfach_En.pdf}
		
		\bibitem{t0_cosmos_de}
		Pascher, Johann (2025). 
		\textit{T0-Modell: Ein vereinheitlichtes, statisches, zyklisches, dunkle-Materie-freies und dunkle-Energie-freies Universum}. 
		T0-Theory Project. 
		\url{https://jpascher.github.io/T0-Time-Mass-Duality/2/pdf/063_cosmic_De.pdf}
		
		\bibitem{t0_cosmos_en}
		Pascher, Johann (2025). 
		\textit{T0-Model: A unified, static, cyclic, dark-matter-free and dark-energy-free universe}. 
		T0-Theory Project. 
		\url{https://jpascher.github.io/T0-Time-Mass-Duality/2/pdf/063_cosmic_En.pdf}
		
		\bibitem{t0_cmb_de}
		Pascher, Johann (2025). 
		\textit{Temperatureinheiten in natürlichen Einheiten: T0-Theorie und statisches Universum}. 
		T0-Theory Project. 
		\url{https://jpascher.github.io/T0-Time-Mass-Duality/2/pdf/061_TempEinheitenCMB_De.pdf}
		
		\bibitem{t0_cmb_en}
		Pascher, Johann (2025). 
		\textit{Temperature Units in Natural Units: T0-Theory and Static Universe}. 
		T0-Theory Project. 
		\url{https://jpascher.github.io/T0-Time-Mass-Duality/2/pdf/061_TempEinheitenCMB_En.pdf}
		
		\bibitem{t0_gravitation_en}
		Pascher, Johann (2025). 
		\textit{Geometric Determination of the Gravitational Constant: From the T0-Model}. 
		T0-Theory Project. 
		\url{https://jpascher.github.io/T0-Time-Mass-Duality/2/pdf/127_gravitationskonstnte_En.pdf}
		
		\bibitem{t0_redshift_de}
		Pascher, Johann (2025). 
		\textit{T0-Theorie: Wellenlängenabhängige Rotverschiebung ohne Distanzannahmen}. 
		T0-Theory Project. 
		\url{https://jpascher.github.io/T0-Time-Mass-Duality/2/pdf/026_T0_Geometrische_Kosmologie_De.pdf}
		
		\bibitem{heisenberg1927}
		Heisenberg, W. (1927). 
		\textit{On the intuitive content of quantum theoretical kinematics and mechanics}. 
		Zeitschrift für Physik, 43(3-4), 172--198.
		
		\bibitem{planck2020}
		Planck Collaboration (2020). 
		\textit{Planck 2018 results. VI. Cosmological parameters}. 
		Astronomy \& Astrophysics, 641, A6. 
		\url{https://doi.org/10.1051/0004-6361/201833910}
		
		\bibitem{codata2018}
		CODATA (2018). 
		\textit{The 2018 CODATA Recommended Values of the Fundamental Physical Constants}. 
		National Institute of Standards and Technology. 
		\url{https://physics.nist.gov/cuu/Constants/}
		
		\bibitem{casimir1948}
		Casimir, H. B. G. (1948). 
		\textit{On the attraction between two perfectly conducting plates}. 
		Proceedings of the Royal Netherlands Academy of Arts and Sciences, 51(7), 793--795.
		
		\bibitem{muon_g2_2021}
		Muon g-2 Collaboration (2021). 
		\textit{Measurement of the Positive Muon Anomalous Magnetic Moment to 0.46 ppm}. 
		Physical Review Letters, 126(14), 141801. 
		\url{https://doi.org/10.1103/PhysRevLett.126.141801}
		
		\bibitem{riess2022}
		Riess, A. G., et al. (2022). 
		\textit{A Comprehensive Measurement of the Local Value of the Hubble Constant with 1 km s$^{-1}$ Mpc$^{-1}$ Uncertainty from the Hubble Space Telescope and the SH0ES Team}. 
		The Astrophysical Journal Letters, 934(1), L7. 
		\url{https://doi.org/10.3847/2041-8213/ac5c5b}
		
		\bibitem{jwst_early}
		Naidu, R. P., et al. (2022). 
		\textit{Two Remarkably Luminous Galaxy Candidates at z $\approx$ 11--13 Revealed by JWST}. 
		The Astrophysical Journal Letters, 940(1), L14. 
		\url{https://doi.org/10.3847/2041-8213/ac9b22}
		
		\bibitem{cobe1992}
		COBE Collaboration (1992). 
		\textit{Structure in the COBE differential microwave radiometer first-year maps}. 
		The Astrophysical Journal Letters, 396, L1--L5. 
		\url{https://doi.org/10.1086/186504}
		
		\bibitem{sparnaay1958}
		Sparnaay, M. J. (1958). 
		\textit{Measurements of attractive forces between flat plates}. 
		Physica, 24(6-10), 751--764. 
		\url{https://doi.org/10.1016/S0031-8914(58)80090-7}
		
		\bibitem{lamoreaux1997}
		Lamoreaux, S. K. (1997). 
		\textit{Demonstration of the Casimir force in the 0.6 to 6 $\mu$m range}. 
		Physical Review Letters, 78(1), 5--8. 
		\url{https://doi.org/10.1103/PhysRevLett.78.5}
		
		\bibitem{einstein1915}
		Einstein, A. (1915). 
		\textit{Die Feldgleichungen der Gravitation}. 
		Sitzungsberichte der Preußischen Akademie der Wissenschaften, 844--847.
		
	\end{thebibliography}
